\documentclass[11pt,a4paper]{article}

\usepackage[T1]{fontenc}
\usepackage[utf8]{inputenc}
\usepackage{lmodern}
\usepackage{amsmath,amssymb}
\usepackage{booktabs}
\usepackage{listings}
\usepackage{xcolor}
\usepackage{geometry}
\usepackage[hidelinks]{hyperref}
\geometry{margin=2.5cm}

\lstdefinelanguage{rewrite}{
  morekeywords={SORT,NUM,CONST,VAR,TERM,FUNC,OTHER,BRA,ATT,ALG,SPEC,ADDSPEC,
                REWR,LIFT,TO,IF,ELSE,FOR,STEP,DO,WHILE,INC,PRT,PRTNUM,OUT,
                DATA,RETURN,QUIT,NIN,IN,LB,LA1,LA2,LA3,LA4,LA5,RA,RU,
                OPT,DEFAULT,RIGHT,COMM,SUBSORT,INTNUM},
  sensitive=true,
  morecomment=[l]{\#},
  morestring=[b]",
}
\lstset{
  language=rewrite,
  basicstyle=\ttfamily\small,
  keywordstyle=\bfseries,
  commentstyle=\itshape\color{black!60},
  columns=fullflexible,
  keepspaces=true,
  frame=single,
  framesep=4pt,
  xleftmargin=6pt,
  xrightmargin=6pt,
  aboveskip=8pt,
  belowskip=8pt,
}

% results of the test-suite measurement (tools/rubitest.py)
\newcommand{\nproblems}{1540}
\newcommand{\ntested}{1539}
\newcommand{\nsolved}{1292}
\newcommand{\psolved}{84.0\%}
\newcommand{\nwrong}{2}
\newcommand{\nfailed}{0}

\title{\texttt{rewrite}: a programming language that learns its\\
       mathematics from rule files}
\author{Rainer Heintzmann\\
        Friedrich-Schiller-Universit\"at Jena, Germany\\
        \texttt{rainer.heintzmann@uni-jena.de}}
\date{}

\begin{document}
\maketitle

\begin{abstract}
\texttt{rewrite} is a small interpreted language for symbolic and numerical
mathematics, written in the early 1990s and revived in 2026. Unlike most
computer algebra systems, its kernel contains no mathematics: data types,
operators, the laws of algebra, differentiation and integration rules are
all defined in plain-text include files and applied by term rewriting. Rule
sets (``algebras'') are first-class objects that can be combined into
sequences and fixpoint iterations. A \texttt{LIFT} statement transfers all
rules of one set of operations to another, so that, for example, the rules
of an abelian group written for $+$ also apply to $\cdot$ in a field.
Numerical evaluation is interleaved with rewriting: as soon as all
arguments of a function are numbers, an attached C function computes the
value. We describe the design of the language, its revival on current
64-bit compilers, and its extension by static subsorts, which add exact
integers and fractions as sorts below the real numbers without changing
any rule of the library, and sorts of non-negative and positive numbers
for which rules like $\sqrt{x^2}=x$ hold. As a case study, the rule-based
integrator Rubi was translated into \texttt{rewrite}. A translator converts
91\% of Rubi's 7422 rules automatically; for this, the language gained ordered rule lists
with a backtracking matcher for optional pattern parts, commutative
operators and condition lists, symbolic helper functions and special
functions. On \nproblems{} integration problems from four independent
textbook collections, \psolved{} of the results are verified correct by
numerical differentiation.
\end{abstract}

\section{Introduction}

Computer algebra systems such as Mathematica~\cite{mathematica} or
SymPy~\cite{meurer2017sympy} come with a large body of built-in mathematical
knowledge. A user can add rules, but the core simplification and the
calculus routines are part of the system and are mostly hidden from the
user. Term rewriting systems such as Maude~\cite{clavel2007maude} take the
opposite view: they provide a general rewriting engine and leave the
mathematics to the user, but they are usually not aimed at numerical work.

\texttt{rewrite} sits between these positions. Its interpreter knows how to
parse expressions with user-defined operators, how to match and replace
terms, and how to call C functions for numerical values. Everything else is
defined in include files: the sort \texttt{real}, the operators \texttt{+},
\texttt{*} and \texttt{\^}, the group, ring and field laws, the
differentiation rules, and numerical algorithms such as Simpson integration.
The language was designed by the author between 1992 and 1996 on Unix and
the Amiga, on top of a scanner and parser framework by Axel T.\ Schreiner.
In 2026 it was ported to current compilers and extended with a test suite,
exact numbers and a rule-based integration library.

The main contributions of this paper are:
\begin{itemize}
  \item a description of a language in which mathematics is a library of
        rewrite rules, including the \texttt{LIFT} mechanism for transferring
        rules between algebraic structures;
  \item the interleaving of symbolic rewriting and numerical evaluation,
        which allows numerical algorithms to be written as rewrite rules;
  \item static subsorts with value conversion, which add exact integers and
        fractions to the language as a library, with the rules of the real
        numbers inherited through \texttt{LIFT};
  \item an automatic translation of the rule-based integrator
        Rubi~\cite{rich2018rubi} into \texttt{rewrite}, with a measurement on
        independent test problems and a discussion of what the translation
        requires from the host system.
\end{itemize}

\section{Language design}

\subsection{Sorts, operators and numbers}

A program first declares sorts (types) and operators. The following
complete program defines real numbers with three operators and evaluates
expressions:

\begin{lstlisting}
SORT real                                 # a data type
NUM real                                  # numbers typed in have sort real
FUNC + : real,real >> real LA2 ATT plus   # left-associative, priority 2
FUNC * : real,real >> real LA3 ATT mal    # binds tighter than +
FUNC ^ : real,real >> real RA  ATT hoch   # right-associative
1+2*3                                     # prints 7
2^3^2                                     # prints 512
CONST x,y : real
x+y*2                                     # prints x + y * 2
\end{lstlisting}

The format words \texttt{LA1}--\texttt{LA5}, \texttt{RA} and \texttt{RU}
define left-associative binary operators of five priorities,
right-associative operators and right unary operators such as the
factorial. \texttt{ATT name} attaches a C function from a table of user
functions. It is called whenever all arguments of the function are numbers;
otherwise the term stays symbolic. \texttt{OTHER} adds further meanings to
an existing name for other argument sorts, so that \texttt{+} can mean
addition of reals, complex numbers or vectors. \texttt{BRA} defines
functions written with brackets, such as a scalar product
$\langle a,b\rangle$.

Each sort has its own representation in C. Real numbers are
\texttt{double} values; the standard library defines complex numbers, real
vectors and matrices, whose values are C arrays. Operations on whole
vectors and matrices are executed by C functions, so they run at compiled
speed.

\subsection{Rules and algebras}

Rules are collected in named rule sets, called \emph{algebras}. A rule has
a left side, a right side, an optional condition and an optional block of
instructions that is executed before the replacement:

\begin{lstlisting}
VAR A : real
ALG simp
SPEC    simp : A*1 = A             # SPEC starts a new rule set
ADDSPEC simp : A+0 = A             # ADDSPEC adds a rule
CONST u : real
REWR (u*1)+0 : simp                # prints u
\end{lstlisting}

\texttt{VAR} declares meta variables that match any term of their sort.
\texttt{NUM} declares meta variables that only match numbers, which makes
rules like ``evaluate if the argument is a number'' easy to write.
Conditions test whether a term occurs in another one (\texttt{IN},
\texttt{NIN}), compare terms in a total order (\texttt{LB}, used for
sorting sums), or evaluate any expression with a boolean result.

Algebras are combined into lists. A list in parentheses is applied
repeatedly until no rule changes the term any more. The standard rule list
of the library is defined like this:

\begin{lstlisting}
SPEC All : Std,Doit,(Linear,Compl,diff,GrRi),AGroup,integr,
           (Subst,GrRi),Std,Zus,unaerweg
\end{lstlisting}

With the library loaded, \texttt{REWR D(sin cos x,x) : All} returns
\texttt{- cos (cos (x ))* sin (x )}.

\subsection{Transferring rules with LIFT}

A distinctive feature of \texttt{rewrite} is the \texttt{LIFT} statement.
It declares that all rules valid for one set of functions and constants
also hold for another set. The standard library writes the laws of a group
once and transfers them step by step:

\begin{lstlisting}
LIFT (Ag+Ag,Ag-Ag,-Ag,0g) TO (Aag+Aag,Aag-Aag,-Aag,0ag)
LIFT (Aag+Aag,Aag-Aag,-Aag,0ag) TO (Ar+Ar,Ar-Ar,-Ar,0r)
LIFT (Aag+Aag,Aag-Aag,-Aag,0ag) TO (Ar*Ar,Ar/Ar,mi(Ar),1r)
\end{lstlisting}

The last line reuses the rules of an abelian group for multiplication,
division, the multiplicative inverse \texttt{mi} and the neutral element 1
of a ring. In the same way, linearity is written once for a generic linear
function and lifted to differentiation, integration, Taylor expansion and
summation. Each symbol can be lifted from one other symbol, so the lifted
operators form a tree, and a rule applies to a term whose operator reaches
the operator of the rule along this tree. This corresponds closely to the
algebraic hierarchies of modern proof assistants~\cite{mathlib2020}, but
here it acts directly on the rewrite rules.

\subsection{Exact numbers as subsorts}

Exact arithmetic was added in 2026 without changing the rules of the
library, by two sorts and two new statements. The sorts \texttt{integer}
(arbitrary size) and \texttt{rational} (fractions) are implemented in C with
the GMP library~\cite{gmp} and declared in an include file:

\begin{lstlisting}
SORT integer
SORT rational
ATT integer  show_int,cmp_int,cpy_int,del_int
ATT rational show_rat,cmp_rat,cpy_rat,del_rat
SUBSORT integer,rational ATT int2rat     # every integer is a fraction
SUBSORT rational,real ATT rat2real       # every fraction is a real number
OTHER + : rational,rational >> rational ATT rat_plus
OTHER / : integer,integer >> rational  ATT int_div
...
LIFT (A+A,A-A,A*A,A/A,-A,mi(A),A^A) TO (Aq+Aq,Aq-Aq,Aq*Aq,Aq/Aq,-Aq,mi(Aq),Aq^Az)
INTNUM natural ATT read_int              # 12 is a natural number, 0.5 real
\end{lstlisting}

\texttt{SUBSORT} declares that every value of the first sort is also a value
of the second, with a C function that converts the value. Sorts remain
static: the sort of a term follows from its parts when it is built, and a
value never changes its sort at run time. A function declared for a wider
sort accepts arguments of its subsorts; if several variants of an operator
fit, the one with the fewest conversions is chosen, and the values are
converted only when the function is calculated. Thus \texttt{1/3+1/6}
uses the rational \texttt{+} and prints \texttt{(1/2)}, \texttt{2\^{}100}
is exact, and \texttt{1/3+0.5} uses the real \texttt{+} and prints
\texttt{0.833333}. The \texttt{LIFT} line lets \texttt{rational} inherit all
rules of \texttt{real}, and a second one does the same for
\texttt{integer}; \texttt{x*1+0*x+(1/2)*x} simplifies to \texttt{(3/2) * x}
with the unchanged rules for real numbers.

\paragraph{Sorts with a known sign.} A sort may have several wider sorts,
and the conversion follows the shortest path in the resulting graph. The
library declares
\[
\mathit{natural} < \mathit{integer} < \mathit{rational} < \mathit{real},\qquad
\mathit{natural} < \mathit{nonneg} < \mathit{real},\qquad
\mathit{positive} < \mathit{nonneg},
\]
where \texttt{nonneg} and \texttt{positive} are real numbers known to be
non-negative or positive. Since the minus sign is an operator, every number
literal is \texttt{natural}, and \texttt{2-5} is the \texttt{integer} $-3$.
Operator variants such as \texttt{nonneg*nonneg >> nonneg} and
\texttt{sqrt : nonneg >> nonneg} propagate the sign through terms. Several
simplification rules are valid only for non-negative or positive
arguments: $\sqrt{t^2}=t$, $(t^a)^b=t^{ab}$, $\ln(pq)=\ln p+\ln q$. They
are now written with meta variables of these sorts. A meta variable whose
sort is a subsort of another only matches terms of that sort or below it,
whereas a variable of sort \texttt{real} still matches any term as in the
original language. Thus \texttt{sqrt(t\^{}2)} simplifies to \texttt{t} for
\texttt{CONST t : nonneg}, but \texttt{sqrt(x\^{}2)} stays as it is for a
real \texttt{x}.

One change in the interpreter made this work. When a rule written for
\texttt{real} rewrites a term with exact numbers, the operators of the new
term are chosen again for the sorts of their actual arguments, from the
leaves upwards and after each numerical step. The original \texttt{LIFT}
mechanism replaced all operators of one sort by those of another; that was
correct for homogeneous structures like groups and rings, but with
subsorts the arguments of one term have different sorts. Re-choosing the
variant repairs this and makes rules for \texttt{real} produce exact
results whenever their inputs are exact. The declarations are a classic
instance of order-sorted algebra~\cite{goguen1992osa}; the novelty here is
that the subsort relation, the conversions and the inheritance of rules
are all defined in the library.

\subsection{Interleaving rewriting and numerics}

After every single rewrite step, the interpreter evaluates all subterms
whose arguments are numbers. Constants that have a value are replaced by
their value when a term is copied. Together with the procedural part of
the language (\texttt{FOR}, \texttt{WHILE}, \texttt{IF}, blocks with local
constants, and user functions), this allows numerical algorithms to be
written as rewrite rules. The library defines Simpson integration as a
rule whose \texttt{DO} block computes the value:

\begin{lstlisting}
ALG rechne
SPEC rechne : I(N,M,A,_x)=_myerg DO { _myerg=0;_zw1=0;_zw2=0  # Simpson
        FOR n=N+_step/2 TO M-_step STEP _step DO { _x=n;_zw1=_zw1+A
        _x=n+_step/2
        _zw2=_zw2+A }
        _x=M-_step/2
        _zw1=_zw1+A
        _x=N;_myerg=A+4*_zw1+2*_zw2
        _x=M;_myerg=(_myerg+A)*_step/2/3
        _x=dummi }
\end{lstlisting}

A second rule set first substitutes the integration variable by the
internal variable \texttt{\_x}. With \texttt{\_step=0.01},
\texttt{REWR I(0,pi,sin(x),x) : numerik,rechne} returns 2. The same
mechanism gives Newton's method and the search for minima and maxima. The
symbolic part can prepare the numerical part: \texttt{newton} expects the
derivative as an argument, which is computed symbolically by \texttt{D}.

The original application of the language was physics. The library contains
physical constants, units and formulas, and an example program computes
Bragg angles, structure factors and X-ray absorption of lithium niobate
crystals from tabulated atomic data.

\section{Case study: rule-based integration}

Differentiation is easy to express with rewrite rules, because every rule
takes the term apart and the product and chain rules always apply.
Integration is harder. Complete decision procedures exist for elementary
functions~\cite{risch1969problem,bronstein2005symbolic}, but they are not
pattern based. Rubi~\cite{rich2018rubi,rubirepo} showed that an extensive
system of carefully ordered rules with conditions can integrate a very
large class of functions. Its 199 rule files contain 7422 rules in
Mathematica syntax, published under the MIT license. We first translated a
small closed subset by hand and then wrote a translator for all of them, to
test how far the mechanisms of \texttt{rewrite} carry.

\subsection{Pattern features}

Rubi's patterns rely on features of the Mathematica pattern matcher that
the original \texttt{rewrite} did not have. A typical rule is

\begin{lstlisting}
Int[(a_. + b_.*x_)^m_, x_Symbol] := (a + b*x)^(m + 1)/(b*(m + 1)) /;
    FreeQ[{a, b, m}, x] && NeQ[m, -1]
\end{lstlisting}

where \texttt{a\_.} matches a missing summand with default 0 and
\texttt{b\_.} a missing factor with default 1. The pattern also matches
$x\,b + a$, because Mathematica matches sums and products in any order, and
the condition combines several tests. We added three corresponding
features to the language: \texttt{OPT a,b} declares meta variables that may
be missing, \texttt{DEFAULT + = 0}, \texttt{DEFAULT * = 1} and
\texttt{DEFAULT RIGHT \^{} = 1} give the values of missing operands,
\texttt{COMM +,*} declares commutative operators, and a rule can have a list
of conditions, all of which must hold. \texttt{NIN(x,a,b,m)} tests that $x$
occurs in none of $a$, $b$, $m$, which is Rubi's \texttt{FreeQ}.

A first implementation expanded every rule into all its variants when it
was defined, so that the matcher stayed unchanged. The number of variants
grows exponentially with the number of optional parameters and commutative
products, however: the median Rubi rule would have about 1900 variants.
Rules defined while one of these features is active are therefore kept in
an ordered \emph{rule list}, as in Mathematica, and matched by a separate
backtracking matcher. Bindings of meta variables are recorded on a trail,
so a failed alternative undoes exactly its own bindings. Swapped operands
and missing optional operands are tried in turn, and the conditions are
checked as part of the match, so a failing condition makes the matcher try
the next alternative and then the next rule. A precomputed set of possible
head symbols per rule skips most rules of a large list without running the
matcher. With these features the rule above translates almost one to one:

\begin{lstlisting}
ADDSPEC RubiRules : IF NIN(Xv,O_a,O_b), (neq(P_m,(-1)))
        Int(((O_a)+((O_b)*(Xv)))^(P_m),Xv) = ...
\end{lstlisting}

Here \texttt{P\_m} is a required and \texttt{O\_a}, \texttt{O\_b} are
optional meta variables. One rule covers $(a+bx)^m$, $(bx+a)^m$, $(a+x)^m$,
$(bx)^m$ and $x^m$, with numeric or symbolic coefficients.

\subsection{Automatic translation}

The translator (900 lines of Python) parses Rubi's rule files with a small
Pratt parser into Mathematica's full form and writes one \texttt{rewrite}
file per Rubi file, in Rubi's loading order and with the original of every
rule as a comment. Pattern variables become meta variables, negative
numbers in patterns become \texttt{NUM} variables with a condition, local
definitions (\texttt{With}, and \texttt{Module} with values only) are
substituted, helper predicates that Rubi defines inside its rule files are
inlined, and a predicate applied to a list, like \texttt{LinearQ[\{u,v\},x]},
becomes a conjunction. As in Mathematica, $e^u$ is a power whose base is a
symbolic constant $e$ of sort \texttt{positive} without a numerical value,
so that the rules for $F^u$ apply to it and $\ln e=1$ stays exact. Some
rules have a
function variable, \texttt{F\_[c\_.*(a\_.+b\_.*x\_)]} with a condition such
as \texttt{TrigQ[F]} or \texttt{MemberQ[\{Sinh,Cosh\},F]}; they are expanded
into one rule per admissible function.

\paragraph{Helper and special functions.} Rubi's conditions and right-hand
sides use about a hundred helper functions, such as \texttt{ExpandIntegrand},
\texttt{PolyQ}, \texttt{BinomialQ}, \texttt{Subst} or
\texttt{ExpandTrigReduce}. They cannot be ordinary \texttt{rewrite}
functions, which are only evaluated when all arguments are numbers. We
added symbolic C functions instead: they receive the argument terms and
return a new term, or nothing if they do not apply. About 60 of them are
implemented, from predicates on the form of a term to polynomial division,
the product-to-sum reduction of trigonometric powers, and Rubi's measures
of simplicity based on the size of a term. Their results are exact when
their inputs are. \texttt{ExpandIntegrand} splits rational functions into
partial fractions, as Mathematica's \texttt{Apart}: with GMP rationals, the
denominator is split into square-free parts, rational roots and the factors
given in the input, and the extended Euclidean algorithm gives the
numerators. Rubi's last general rule, $\int u\,dx=\int
\texttt{ExpandIntegrand}[u]\,dx$ if the result is a sum, then integrates
every rational function term by term. A condition like
\texttt{SumQ[ExpandIntegrand[u,x]]} needs its inner helper evaluated
first; \texttt{rewrite} calculated conditions from the outside in, so the
rule never applied until this was changed. Rubi's substitutions rest on
\texttt{FunctionOfQ} and \texttt{SubstFor}, which decide whether the
integrand is a function of, say, $\sin w$ and write it in a new variable,
also through identities such as $\cos^2w=1-\sin^2w$; with these two C
functions, 34 substitution rules for trigonometric and hyperbolic functions
became translatable. \texttt{DerivativeDivides[y,u,x]} returns $u/y'$ if it
is free of $x$, and \texttt{False} otherwise; it is a C function that stays
unevaluated when it does not apply, and \texttt{FalseQ} tests for exactly
that, so that 20 rules of the form $\int u\,f(y)\,dx=q\int f(t)\,dt$ with
$t=y$ could be translated. The substitution $x\to F^{gx}$ for functions of
exponentials is a C function that takes the place of Rubi's rule, which
uses pattern tests the translator cannot express. Many results of Rubi need special functions; 27 of them,
from the error function and the Fresnel and exponential integrals to the
incomplete Gamma function, elliptic integrals in Carlson's
form~\cite{carlson1995} and the hypergeometric function ${}_2F_1$, were
added as C functions with Mathematica's conventions, each with a derivative
rule. Their values agree with mpmath~\cite{mpmath} to within $10^{-11}$.

\paragraph{The host system's simplification.} Rubi relies on Mathematica's
automatic simplification, which \texttt{rewrite} does not have. Most of the
work after the first translation went into reproducing the parts Rubi
depends on. A normal-form rule set writes $x-y$ as $x+(-1)\,y$, $x/y$ as
$x\,y^{-1}$, $\sqrt{x}$ as $x^{1/2}$ and $(uv)^k$ as $u^kv^k$ for integer
$k$ (in the translated patterns as well), but, unlike our first version and
like Mathematica, it does not multiply out products with sums: Rubi's rules
for $(c+d\tan x)/(a+b\tan x)$ expect the product. Polynomials in the
integrand are brought into the form $c_0+c_1x+c_2x^2+\dots$ with collected
coefficients and their numeric content in front, and every rule result is
simplified trivially ($u^0$, $u^1$, $1\cdot u$, $0+u$). In Mathematica,
some rules that make no progress are harmless because exact arithmetic or
other rules prevent a loop; in \texttt{rewrite} they looped. The rule list
therefore skips a rule whose result equals the term, contains a NaN or a
division by zero, or closes a cycle such as expanding and factoring again
(the result was replaced before, and the chain of the last replacements,
also of subterms, leads back to the term), and a rule that keeps rewriting
the same term. Only the chain counts: a term that was replaced in another
branch, like $\int b\,dx$, may come up again. The matcher also had to learn that an optional
variable already bound to its default value may be missing a second time,
as in $(c+d\tan(e+fx))/(a+b\tan(e+fx))$ with $e=0$, and that products are
flat: Mathematica's pattern $u\,(a+b\arctan x)$ matches
$x^2\arctan x/(1+x^2)$ with $u=x^2/(1+x^2)$, so for a commutative operator
the matcher also tries every split of up to six nested operands into two
groups. Finally, Rubi's
substitutions $\int u\,dx$ with $x\to\sqrt{a+bx}$ rely on Mathematica
treating the new variable as non-negative. \texttt{Subst[Int[u,x],x,v]}
therefore integrates in a fresh constant of sort \texttt{nonneg} when $v$ is
a square root or otherwise non-negative, so that $\sqrt{x^2}=x$ holds
inside; with this, $\int\sin\sqrt{1+x}/\sqrt{1+x}\,dx=-2\cos\sqrt{1+x}$.
Mathematica also writes $(b\tan u)^3$ as $b^3\tan^3u$, so a pattern
$(b\,\tan u)^n$ with integer $n$ always sees $b=1$ there, and some rules
silently drop $b$. The translator adds the condition $b=1$ (or $b=0$ in a
sum) to every rule whose optional variable appears neither in a condition
nor in the result.

Two parts of Rubi lie outside its rule files. Most trigonometric rules exist
only for sine, secant and tangent; Rubi's utility function
\texttt{UnifyInertTrigFunction} first rewrites cosine forms, such as
$a+b\cos(e+fx)$ as $a+b\sin(e+\pi/2+fx)$. Its 69 definitions are translated
by the same program into rules that run before the trigonometric ones, and
the output rules turn $\sin(\pi/2+u)$ back into $\cos u$. For hyperbolic
functions, Rubi uses $\sinh u=-i\sin(iu)$, which needs complex numbers.
Instead, a last rule writes hyperbolic functions for which no other rule
applies with exponentials, $\sinh u=(e^u-e^{-u})/2$, so that
$\int x\sinh x\,dx$ and $\int\cosh^2x\,dx$ are found.

The complete integration is the algebra list
\begin{lstlisting}
SPEC RubiIntegrate : (RubiNorm,IntRubi,RubiRules),RubiCollect,IntOut,All
\end{lstlisting}
where \texttt{RubiCollect} combines like summands of the result and
\texttt{IntOut} restores the usual output form.

\subsection{Results}

The translator converts 6830 of Rubi's 7422 rules (92\%), which become
7321 \texttt{rewrite} rules after the expansion of function variables. The
largest group of the remaining rules needs complex numbers (245 rules);
others use the rest of Rubi's substitution machinery
(\texttt{FunctionOfExponentialQ}), general pattern tests
(\texttt{MatchQ}), function variables without a finite set of functions,
or the generalized hypergeometric function. Loading all translated rules takes
about half a second, and a typical integral a fraction of a second.
Table~\ref{tab:int} shows examples.

\begin{table}[ht]
\centering
\small
\begin{tabular}{@{}ll@{}}
\toprule
Integrand & Result of \texttt{REWR Int(f,x) : RubiIntegrate} \\
\midrule
\texttt{1/(1+x\^{}3)} & \texttt{(1/3)*ln(1+x) + 0.57735*arctan(-0.57735+1.1547*x)} \\
                      & \texttt{\ \ - (1/6)*ln(1+x\^{}2-x)} \\
\texttt{x/(1+x\^{}4)}  & \texttt{(1/2)*arctan(x\^{}2)} \\
\texttt{sqrt(1-x\^{}2)} & \texttt{(1/2)*arcsin(x) + (1/2)*(1-x\^{}2)\^{}(1/2)*x} \\
\texttt{x*sqrt(1+x)}  & \texttt{(2/5)*(1+x)\^{}(5/2) - (2/3)*(1+x)\^{}(3/2)} \\
\texttt{sin(x)\^{}2*cos(x)\^{}3} & \texttt{(1/3)*sin(x)\^{}3 - (1/5)*sin(x)\^{}5} \\
\texttt{1/(1+tan(x))\^{}2} & \texttt{(1/2)*ln(cos(x)+sin(x)) - (1/2)/(1+tan(x))} \\
\texttt{x*arcsec(x)} & \texttt{(1/2)*x\^{}2*arcsec(x) - (1/2)*(1-1/x\^{}2)\^{}(1/2)*x} \\
\texttt{exp(-x\^{}2)} & \texttt{0.886227*erf(x)} \\
\texttt{sin(x)/x}     & \texttt{sinint(x)} \\
\texttt{1/sqrt(1-x\^{}4)} & \texttt{ellipticf(arcsin(x),-1)} \\
\texttt{(1+x\^{}3)\^{}(1/3)} & \texttt{hyp2f1((-1/3),(1/3),(4/3),-x\^{}3)*x} \\
\bottomrule
\end{tabular}
\caption{Indefinite integrals computed with the automatically translated
Rubi rules (output slightly condensed). Rational coefficients are exact;
irrational ones such as $1/\sqrt{3}$ and $\sqrt{\pi}/2$ are real numbers.}
\label{tab:int}
\end{table}

To measure the translation, a test harness runs problems of Rubi's own test
suite~\cite{rubitests} through \texttt{RubiIntegrate}. It does not compare
the result with Rubi's antiderivative, which may differ by a constant or in
form, but checks it numerically: the derivative of the result, computed by
a five-point difference quotient, must equal the integrand at several
points where the integrand is real, with fixed values for symbolic
parameters. The four collections of independent problems in the test
suite (from the textbooks of Apostol, Hearn, Stewart and Timofeev,
\nproblems{} problems, of which \ntested{} use only functions that
\texttt{rewrite} has) gave the results in Table~\ref{tab:results}.

\begin{table}[ht]
\centering
\small
\begin{tabular}{@{}lrr@{}}
\toprule
Result & Problems & Share \\
\midrule
solved and verified       & \nsolved & \psolved \\
no antiderivative found   & 229 & 14.9\% \\
wrong result              & \nwrong  & 0.1\% \\
crash or time-out         & \nfailed & 0.0\% \\
not verifiable            & 16 & 1.0\% \\
\bottomrule
\end{tabular}
\caption{Results on the \ntested{} tested problems of four independent
integration problem collections in Rubi's test suite.}
\label{tab:results}
\end{table}

The development of the translation went through several such measurements;
each round of fixes, such as the normal form without multiplying out, the
matching of optional variables bound to their default, the special
functions and the missing helpers, raised the share of verified results.
During a first measurement, 41\% of the problems were solved, 55\% before
the cosine unification, the hyperbolic fallback and the symbolic $e$, 62\%
after them, 79\% after partial fractions, the inside-out evaluation of
nested helpers and the flat matching of products, and 84\% with Rubi's
substitutions. The two
results counted as wrong contain special functions evaluated outside their
real domain (${}_2F_1$ at an argument $e^{7x}/3>1$, the incomplete Gamma
function of a negative argument), where the check compares different
complex branches. A run takes about three minutes.

The translation also served as a test of the library itself. It revealed
two incorrect simplification rules in the standard algebra, which turned
$\ln(x+y)$ into $\ln x\cdot\ln y$ and $a^{b^c}$ into $a^{bc}$, two incorrect
derivatives (of $\arcsin$ and $\arccos$), and three rules that are only
valid for positive bases: $(ab)^c=a^c b^c$, $(a^b)^c=a^{bc}$, which turned
$\sqrt{x^2}$ into $x$, and $\ln(ab)=\ln a+\ln b$. They now apply only for
integer exponents or a positive numerical factor, and for terms of the
sorts \texttt{nonneg} and \texttt{positive}. The numerical check of the
harness found these errors, because Rubi's results were differentiated at
negative arguments.

\section{Implementation and revival}

The interpreter consists of about 8800 lines of C, of which about 3900 were
added in 2026 for the rule lists, subsorts, exact numbers, special and
helper functions, and a grammar of about 490 lines for the scanner and
parser generators flex and bison. Terms are binary trees of nodes taken
from memory pools; symbols carry their sort, value and the links used by
\texttt{LIFT}. The rules of an algebra are merged into a single tree, in
which rules with a common left-hand-side structure share nodes and
alternatives branch off where they differ; rule lists are kept in order
instead. Rules are matched against every subterm, innermost first, and a
single \texttt{REWR} stops after $10^5$ steps, so that a looping rule set
cannot hang the interpreter.

The program was written for 32-bit Unix and the Amiga and no longer worked
with current tools. The port to 64-bit Linux and Windows (MSYS2) required
only small changes to the C code itself. The main problems were changed
semantics in the scanner generator, the release of static memory that old
C libraries tolerated, and errors in the handling of declarations inside
loops and included files. A regression test suite with 28 test programs
now covers the core language, the library, the pattern features, exact
numbers, helper and special functions, the integration rules and the
original example programs; the integration tests check every result by
differentiating it again. Memory leaks were removed with an instrumented
build that counts allocations per source line, so that repeated
calculations run in constant memory. All parser structures are now
accessed through named macros instead of chains like
\texttt{np->r->r->l}; for most source files the compiled code stayed
byte-identical, which confirmed that this cleanup did not change the
behaviour. \texttt{rewrite} runs interactively with a prompt, runs program
files given as arguments, or reads a program from a pipe for scripts and
the tests.

Performance is adequate for interactive use. Loading the standard library
and computing a derivative takes well below 0.1\,s on a current laptop;
loading the 7321 integration rules about half a second. A scalar loop
written in the language itself needs about 0.25\,$\mu$s per iteration
(measured for $10^7$ iterations of \texttt{s = s + k*k}). We estimate this
to be about two orders of magnitude slower than an equivalent compiled
loop. Operations on whole vectors and matrices run in C.

\section{Discussion}

\paragraph{Strengths.} The main strength of \texttt{rewrite} is that all
mathematics is visible and changeable. Correcting the derivative of
$\arcsin$ meant editing one line of a text file, not the interpreter.
\texttt{LIFT} allows algebraic structure to be expressed once and reused;
with subsorts, it also turned out to be the right mechanism for a numeric
tower, since exact fractions inherited all rules of the real numbers
without any change to them. \texttt{NUM} variables together with conditions
give precise control over when a rule fires. The interleaving of rewriting
and numerical evaluation makes it natural to derive formulas symbolically
and evaluate them numerically in the same program, for example derivatives
for Newton's method.

\paragraph{Limitations.} Commutative matching only swaps the two operands
of one binary operator. Associativity is handled by normalization, so the
integrand must be brought into a fixed form before the rules are tried.
Complex numbers exist as a sort of the library, but not yet as part of the
number tower, so rules that need the imaginary unit are not translated,
and logarithms of negative arguments evaluate to NaN although their
derivatives are correct. The sign of a term is only known through the
sorts of its operator variants, so for example $e^x$ is not known to be
positive. Nothing checks that a rule set
terminates or is confluent~\cite{baader1998term}, and the order of rules
matters; the safeguards against loops in rule lists are heuristics.
Numeric conditions such as $m\ne-1$ do not hold for symbolic $m$, which is
more conservative than Rubi. Error messages are minimal.

\paragraph{Related work.} Rule-based simplification is at the core of
Mathematica~\cite{mathematica} and of Rubi~\cite{rich2018rubi}. General
rewriting logic, with order-sorted signatures, is provided by
Maude~\cite{clavel2007maude,goguen1992osa}. More recently, equality
saturation with e-graphs, as in egg~\cite{willsey2021egg} and
Metatheory.jl~\cite{cheli2021metatheory}, avoids the dependence on rule
order by representing many equivalent terms at once. Proof assistants
organize mathematical knowledge in algebraic hierarchies and simplification
rule sets~\cite{mathlib2020}. \texttt{rewrite} anticipated several of these
ideas in a small and self-contained system.

\paragraph{Outlook.} Natural next steps are complex numbers on top of the
number tower, which would also open the 245 Rubi rules with the imaginary
unit but require complex versions of the helper and special functions;
Rubi's substitution machinery; an associative-commutative matcher for sums
and products of any length; and code generation that turns symbolically
derived expressions into compiled numerical code.

\section{Conclusion}

\texttt{rewrite} shows that a small interpreter can host a complete
computer algebra library when the mathematics is expressed as rewrite rules
and rule sets can be combined and lifted between algebraic structures and
number sorts. The automatic translation of 91\% of Rubi demonstrates that
established rule-based algorithms can be imported, and it clarifies which
matching features and which implicit simplifications of the host system
such an import requires.

\section*{Availability}

The source code, the standard library, the integration rules, the
translator, the test harness and the test suite are available at
\url{https://github.com/RainerHeintzmann/rewrite} under the MIT
license~\cite{rewriterepo}. The translated integration rules are derived
from Rubi~\cite{rubirepo} and keep its MIT license notice.

\section*{Acknowledgements}

The scanner and parser framework of the original program is based on code
by Axel T.\ Schreiner. The 2026 port, the test suite, the memory leak
analysis, the exact numbers, the special functions and the translation of
the integration rules were carried out with the assistance of an AI coding
assistant (Claude, Anthropic).

\bibliographystyle{plain}
\bibliography{references}

\end{document}
